\section{Step 25: Admit scratch storage}
The record validator needs a working frame even though the indexer produces
o displayed pixels. Read context offset sixteen and reject a capacity below
16\,384. Read the scratch address at offset eight and require its low two bits
to be zero. Preserve that address at stack offset 80 for nested record calls.

Fill \texttt{index\_context}. Neither reading a pointer nor checking alignment
proves that an allocation exists; the caller supplies the readable and writable
windows promised by this interface. Reject an invalid scratch declaration before
reading file bytes. No context output is stored during admission.

\paragraph{Observe} Prepare \texttt{s25-1} and stop before
\texttt{index\_context\_ready}. At \texttt{index\_context\_ready}, the native doubleword
at \texttt{\$sp+80} is the address four bytes after
\texttt{index\_scratch}. The complete context still equals its entry state.

\Needspace{7\baselineskip}
\paragraph{Check} What may a rejected late record change?
\begin{itemize}
  \item[A.] The public directory.
  \item[B.] The published frame count.
  \item[C.] The active scratch image.
\end{itemize}
